%%%%%%%%%%%%%%%%%%%%%%%%%%%%%%%%%%%%%%%%%%%%%%%%%%%%%%%%%%%%%%%%%%%%%%%%%%%
%
% Generic template for TFC/TFM/TFG/Tesis
%
% By:
%  + Javier Macías-Guarasa.
%    Departamento de Electrónica
%    Universidad de Alcalá
%  + Roberto Barra-Chicote.
%    Departamento de Ingeniería Electrónica
%    Universidad Politécnica de Madrid
%
% By: + Javier Macías-Guarasa. Departamento de Electrónica Universidad de Alcalá + Roberto Barra-Chicote. Departamento de Ingeniería Electrónica Universidad Politécnica de Madrid
%
% Based on original sources by Roberto Barra, Manuel Ocaña, Jesús Nuevo, Pedro Revenga, Fernando Herránz and Noelia Hernández. Thanks a lot to all of them, and to the many anonymous contributors found (thanks to google) that provided help in setting all this up.
%
% See also the additionalContributors.txt file to check the name of additional contributors to this work.
%
% If you think you can add pieces of relevant/useful examples, improvements, please contact us at (macias@depeca.uah.es)
%
% You can freely use this template and please contribute with comments or suggestions!!!
%
%%%%%%%%%%%%%%%%%%%%%%%%%%%%%%%%%%%%%%%%%%%%%%%%%%%%%%%%%%%%%%%%%%%%%%%%%%%

\chapter{Estructura del documento}
\label{cha:estructura-documento}

En este capítulo recorremos las partes que forman el documento principal para que sepas dónde escribir cada contenido. Las describimos todas para que tengas una referencia completa, pero no tendrás que modificar muchas de ellas porque la plantilla las genera o selecciona automáticamente.

\section{Estructura del documento generado por la plantilla}
\label{sec:estr-del-docum}

La estructura del documento principal se define en \texttt{Book/book.tex}. En este apartado describimos todos sus elementos para que puedas utilizarlo como referencia cuando necesites localizar una parte concreta del documento. No tendrás que editar ni configurar todos ellos: algunos se generan automáticamente y sus órdenes de inclusión deben mantenerse sin cambios.

El documento se divide en tres partes gestionadas por las órdenes \texttt{\textbackslash{}frontmatter}, \texttt{\textbackslash{}mainmatter} y \texttt{\textbackslash{}backmatter}:

\begin{itemize}
  \item \texttt{\textbackslash{}frontmatter}: contiene las cubiertas, los documentos preliminares, la dedicatoria, los agradecimientos, los resúmenes y los distintos índices y listados.
  \item \texttt{\textbackslash{}mainmatter}: contiene los capítulos principales, los apartados técnicos opcionales, la bibliografía y los apéndices.
  \item \texttt{\textbackslash{}backmatter}: contiene la contraportada cuando corresponde al tipo de trabajo seleccionado.
\end{itemize}

Estas órdenes controlan la numeración y la organización interna del libro. No necesitas modificarlas ni cambiar su posición.

La estructura distribuida con la plantilla incluye los siguientes elementos:

\begin{itemize}
  \item Portada y páginas iniciales, seleccionadas automáticamente por \texttt{cover/cover.tex} a partir de la titulación configurada mediante \texttt{\textbackslash{}myDegree}. Normalmente solo tendrás que configurar los datos de \texttt{Config/myconfig.tex}, sin modificar \texttt{cover/cover.tex} ni las plantillas institucionales del directorio \texttt{cover/}.
  \item Documentos PDF preliminares. Inclúyelos únicamente cuando los exija la normativa de tu trabajo mediante \texttt{\textbackslash{}includepdf}. El fichero \texttt{letters/sampleLetter.pdf} muestra cómo incorporar todas las páginas de un documento, y \texttt{letters/sampleLetter-pages.pdf} permite probar la selección de páginas concretas.
  \item Dedicatoria, incluida desde \texttt{dedication/dedicatoria.tex}. Puedes editarla o comentar su orden \texttt{\textbackslash{}input} si no quieres utilizarla.
  \item Agradecimientos, incluidos desde \texttt{acknowledgements/agradecimientos.tex}. También puedes editarlos u omitirlos.
  \item Definiciones de acrónimos, incluidas desde \texttt{acronyms/defacronymsgl.tex}, y listado correspondiente, configurado en \texttt{acronyms/acronymsgl.tex}. Puedes comentar ambos elementos cuando no utilices acrónimos; la sección~\ref{sec:uso-de-acronimos} explica su definición y uso.
  \item Definiciones de símbolos, incluidas desde \texttt{symbols/defsymbolsgl.tex}, y listado correspondiente, configurado en \texttt{symbols/symbolsgl.tex}. Puedes comentar ambos elementos cuando no utilices símbolos; la sección~\ref{sec:simbolos} explica su definición y uso.
  \item Resumen en español, incluido desde \texttt{abstract/resumen.tex}, y resumen en inglés, incluido desde \texttt{abstract/abstract.tex}. Normalmente tendrás que editar ambos ficheros con el contenido de tu trabajo.
  \item Resumen extendido en español, disponible en \texttt{abstract/resumen-extendido.tex}. Su inclusión está comentada de forma predeterminada; actívala únicamente cuando lo requiera tu tipo de trabajo.
  \item Índice general, índice de figuras e índice de tablas, generados automáticamente desde \texttt{cover/toc+lof+lot.tex}. No necesitas editar este fichero ni los ficheros auxiliares que genera \LaTeX{} para construir los índices.
  \item Listados adicionales, configurados en \texttt{cover/extralistings.tex}. El ejemplo incluye listados para fragmentos de código, algoritmos y vídeos; la sección~\ref{sec:intro-videolink} explica el ejemplo dedicado a estos últimos. Puedes editar o comentar esta configuración si no la necesitas.
  \item Capítulos principales, incluidos desde el directorio \texttt{chapters/} mediante \texttt{chapters/contents.tex}. El documento distribuido contiene esta guía, un capítulo de elementos básicos que incluye ejemplos de figuras y tablas, y otro de composición avanzada. Cuando prepares tu trabajo, sustituye esta selección por tus propios capítulos o utiliza el esqueleto académico disponible en \texttt{chapters/bare/}.
  \item Pliego de condiciones, disponible en \texttt{pliego/pliego.tex}, y presupuesto, disponible en \texttt{presupuesto/presupuesto.tex}. Sus órdenes de inclusión están comentadas de forma predeterminada; actívalas únicamente cuando necesites estos documentos. Los ejemplos proceden de un trabajo fin de carrera de Jesús Martínez y se conservan como material de referencia; debido a su antigüedad, contrástalos con la normativa y los formatos actuales.
  \item Bibliografía, impresa mediante \texttt{biblio/bibliography.tex}. Añade tus referencias a los ficheros \texttt{.bib} declarados en \texttt{biblio/bibliofiles.tex} mediante las variables \texttt{\textbackslash{}mybibfileOne}, \texttt{\textbackslash{}mybibfileTwo} y las sucesivas variables disponibles; normalmente no necesitarás modificar \texttt{biblio/bibliography.tex}.
  \item Apéndices, agrupados mediante el entorno \texttt{appendices} e incluidos desde \texttt{appendix/contents.tex}. El documento distribuido dedica un apéndice al código fuente, los algoritmos y los listados adicionales. Puedes reemplazarlos, añadir otros nuevos u omitirlos según las necesidades de tu trabajo.
  \item Contraportada, seleccionada automáticamente mediante \texttt{cover/backpage.tex} cuando existe una plantilla aplicable al tipo de trabajo y a la titulación configurados. Normalmente no necesitarás modificar este fichero ni la plantilla institucional seleccionada.
\end{itemize}

En la práctica, normalmente trabajarás en los capítulos, los resúmenes, la bibliografía y, si los utilizas, los acrónimos, los símbolos y los apéndices. Modifica las órdenes de inclusión de \texttt{Book/book.tex} solo cuando necesites añadir u omitir alguno de estos contenidos; conserva sin cambios las órdenes estructurales y los elementos seleccionados automáticamente. Cuando incorpores un PDF mediante \texttt{\textbackslash{}includepdf}, mantén la orden \texttt{\textbackslash{}clearemptydoublepage} posterior para que el siguiente elemento comience en la página adecuada.


%%% Local Variables:
%%% TeX-master: "../book"
%%% End:
